\documentclass[11pt,a4paper]{article}
\usepackage[utf8]{inputenc}
\usepackage[T1]{fontenc}
\usepackage[spanish,mexico]{babel}
\decimalpoint
\usepackage{geometry}
\geometry{margin=2.5cm, top=3cm, bottom=3cm}
\usepackage{lmodern}
\renewcommand{\familydefault}{\sfdefault}
\usepackage{microtype}

\usepackage{amsmath,amssymb}
\usepackage{booktabs,tabularx,array,multirow,longtable}
\usepackage{graphicx,float,caption}
\captionsetup{font=small, labelfont=bf}
\usepackage{enumitem,parskip}
\usepackage{fancyhdr,titlesec}
\usepackage[most,breakable]{tcolorbox}
\usepackage{hyperref}
\usepackage{listings}

%% Paleta UAM-A
\definecolor{uamred}{RGB}{200,45,35}
\definecolor{uamlightred}{RGB}{248,232,232}
\definecolor{uamblue}{RGB}{28,28,28}
\definecolor{uamgray}{RGB}{100,100,100}
\definecolor{okgreen}{RGB}{56,142,60}
\definecolor{codebg}{RGB}{250,250,250}

%% Secciones
\titleformat{\section}
  {\large\bfseries\color{uamblue}}
  {\textcolor{uamred}{\thesection}}{0.8em}{}
  [\vspace{3pt}\color{uamred}\hrule height 0.5pt\vspace{2pt}]
\titleformat{\subsection}
  {\normalsize\bfseries\color{uamblue}}{\thesubsection}{0.8em}{}
\titleformat{\subsubsection}
  {\normalsize\itshape\color{uamgray}}{\thesubsubsection}{0.8em}{}

%% Encabezado/pie
\pagestyle{fancy}\fancyhf{}
\fancyhead[L]{\small\color{uamgray}Modificaciones 2026 --- División de Ciencias Básicas e Ingeniería · UAM-A}
\fancyhead[R]{\small\color{uamgray}Documentación Técnica}
\fancyfoot[C]{\small\thepage}
\renewcommand{\headrulewidth}{0.4pt}
\setlength{\headheight}{15pt}

%% Cajas
\tcbuselibrary{skins,breakable}
\newtcolorbox{notabox}{enhanced,breakable,
  colback=uamlightred,colframe=uamlightred,
  borderline west={3pt}{0pt}{uamred},
  sharp corners,left=10pt,right=8pt,top=5pt,bottom=5pt,
  before upper={\textbf{\color{uamred}Nota}\enspace}}
\newtcolorbox{warnbox}{enhanced,breakable,
  colback=black!4,colframe=black!4,
  borderline west={3pt}{0pt}{uamgray},
  sharp corners,left=10pt,right=8pt,top=5pt,bottom=5pt,
  before upper={\textbf{Advertencia}\enspace}}
\newtcolorbox{okbox}{enhanced,breakable,
  colback=green!6,colframe=green!6,
  borderline west={3pt}{0pt}{okgreen},
  sharp corners,left=10pt,right=8pt,top=5pt,bottom=5pt}

%% Código fuente
\lstset{
  basicstyle=\ttfamily\small,
  backgroundcolor=\color{codebg},
  frame=single,
  framerule=0pt,
  rulecolor=\color{uamred!30},
  xleftmargin=1em,
  xrightmargin=1em,
  breaklines=true,
  breakatwhitespace=true,
  showstringspaces=false,
  inputencoding=utf8,
}

\hypersetup{colorlinks=true, linkcolor=uamblue, urlcolor=uamred}

%% Macros locales
\newcommand{\RES}{\textsc{res}}
\newcommand{\IA}{\textsc{ia}}

%%============================================================
\begin{document}
%%============================================================

%% ── Portada ───────────────────────────────────────────────
\begin{titlepage}
\centering
\includegraphics[width=0.54\textwidth,trim={48bp 64bp 473bp 35bp},clip]%
  {/Users/nowhere/Claude_Projects/Personal/logo_UAM.png}\\[0.6em]
{\color{uamgray} División de Ciencias Básicas e Ingeniería}\\[0.3em]

\vfill
{\color{uamred}\rule{\textwidth}{1pt}}\\[0.8em]
{\fontsize{22}{26}\selectfont
Sistema de Revisión Automatizada\\[0.3em]
de Planes y Programas de Estudio}\\[0.5em]
{\large\color{uamgray}
Documentación técnica de agentes y documentos de salida\\
DCBI --- UAM-Azcapotzalco --- 2026}\\[0.3em]
{\color{uamred}\rule{\textwidth}{1pt}}

\vfill
\begin{tabular}{rl}
  \textbf{Responsable:} & Dr.\ César Simón López-Monsalvo \\
  \textbf{Departamento:} & Ciencias Básicas \\
  \textbf{Proceso:} & Actualización curricular 2026 \\
  \textbf{Versión:} & 1.0 \\
  \textbf{Fecha:} & junio de 2026 \\
\end{tabular}
\end{titlepage}

\tableofcontents
\newpage

%%============================================================
\section{Propósito del sistema}
%%============================================================

Este sistema automatiza la revisión de las propuestas de modificación a planes y
programas de estudio de las diez licenciaturas de la División de Ciencias Básicas e
Ingeniería (DCBI) de la UAM-Azcapotzalco, correspondientes al proceso de actualización
curricular 2026.

El objetivo es verificar que cada propuesta cumpla con los requisitos establecidos en el
Reglamento de Estudios Superiores (\RES{}) de la UAM, en particular los artículos 34, 44,
45 y 56, y que la información contenida en los documentos entregados por los coordinadores
sea coherente e íntegra. Los resultados se exportan a archivos Excel en el formato del
\textit{Formato de Revisión de Planes y Programas} institucional.

%%──────────────────────────────────────────────────────────
\subsection{Licenciaturas cubiertas}
%%──────────────────────────────────────────────────────────

\begin{center}
\begin{tabular}{clp{6cm}}
\toprule
\textbf{\#} & \textbf{Licenciatura} & \textbf{Carpeta de entrada} \\
\midrule
1  & Ing. Ambiental              & \texttt{Ing.\ Ambiental/} \\
2  & Ing. Civil                  & \texttt{Ing.\ Civil/} \\
3  & Ing. Computación            & \texttt{Ing.\ Computacion/} \\
4  & Ing. Eléctrica (LIETS)      & \texttt{Ing.\ Electrica/Ing.\ Electrica/} \\
5  & Ing. Electrónica            & \texttt{Ing.\ Electronica/} \\
6  & Ing. Física                 & \texttt{Ing.\ Fisica/} \\
7  & Ing. Industrial             & \texttt{Ing.\ Industrial/Ing.\ Industrial/} \\
8  & Ing. Mecánica               & \texttt{Ing.\ Mecanica/Ing.\ Mecanica/} \\
9  & Ing. Metalúrgica            & \texttt{Ing.\ Metalurgica/Ing.\ Metalurgica/} \\
10 & Ing. Química                & \texttt{Ing.\ Quimica/} \\
\bottomrule
\end{tabular}
\end{center}

%%============================================================
\section{Arquitectura general}
%%============================================================

El sistema se ejecuta desde \texttt{revisor/main\_v2.py} e instancia una cadena de agentes
especializados que procesan los documentos de cada licenciatura en secuencia.
El flujo completo es el siguiente:

\begin{lstlisting}[language={}]
main_v2.py
  |
  +-- Agente 1: inventario.py
  |     `-- Por cada licenciatura: localiza y clasifica todos los documentos
  |
  +-- Agente 2: extractor_plan.py  (pertinencia y plan de estudios)
  |     `-- Extrae texto completo y tabla de UEAs de .docx / .pdf
  |
  +-- Agente 3: revisor_ia.py  (Claude API -- claude-sonnet-4-6)
  |     +-- Seccion I:  analiza pertinencia contra Art. 34 RES
  |     `-- Seccion II: analiza plan de estudios contra Art. 44 RES
  |
  +-- Agente 4: extractor_uea.py
  |     `-- Por cada UEA: extrae encabezado y verifica secciones requeridas
  |
  +-- Agente 5: correspondencia_v2.py
  |     `-- Cruza tabla de UEAs del plan contra archivos de programa
  |
  `-- Agente Excel: excel_v2.py
        +-- Revision_DCBI_2026_Planes_<fecha>.xlsx    (Secciones I y II)
        `-- Revision_DCBI_2026_Programas_<fecha>.xlsx (Secciones III y IV)
\end{lstlisting}

Los agentes 2--5 se ejecutan por licenciatura; el Agente Excel consolida la totalidad de
los resultados al final.

%%============================================================
\section{Descripción de los agentes}
%%============================================================

%%──────────────────────────────────────────────────────────
\subsection{Agente~1 --- Inventario (\texttt{inventario.py})}
%%──────────────────────────────────────────────────────────

\subsubsection{Función}

Recorre de forma recursiva la carpeta de cada licenciatura y clasifica cada archivo
encontrado según el tipo de documento que representa. Esta clasificación alimenta a
todos los demás agentes: sin un inventario correcto, los agentes de extracción no saben
qué archivo procesar.

\subsubsection{Lógica de clasificación}

La clasificación se realiza en dos pasos:

\begin{enumerate}[nosep]
  \item \textbf{Nombre del archivo}: se buscan palabras clave en el nombre del archivo
        (sin extensión, normalizado a minúsculas con guiones y subrayados reemplazados
        por espacios).
  \item \textbf{Nombre de la carpeta padre}: si el nombre del archivo no contiene las
        palabras clave, se revisa el nombre de la carpeta inmediatamente superior.
        Este segundo paso es necesario para casos como Ing.\ Computación, cuyo
        documento de pertinencia tiene un nombre genérico pero reside en la carpeta
        \texttt{1.\ Documento de pertinencia y justificacion de la modificacion del PE}.
\end{enumerate}

\begin{center}
\begin{tabular}{lp{8cm}}
\toprule
\textbf{Tipo} & \textbf{Palabras clave} \\
\midrule
\texttt{pertinencia}      & \texttt{pertinencia}, \texttt{justificaci} \\
\texttt{plan\_estudios}   & \texttt{plan de estudio}, \texttt{plan\_estudio}, \texttt{plan estudio} \\
\texttt{mapa\_curricular} & \texttt{mapa}, \texttt{malla}, \texttt{curricular}, \texttt{boligrama} \\
\texttt{equivalencias}    & \texttt{equivalencia} \\
\texttt{oficio}           & \texttt{oficio} \\
\texttt{otro}             & (no coincide con ninguna categoría) \\
\bottomrule
\end{tabular}
\end{center}

\begin{notabox}
Los archivos con la cadena \texttt{anexo} en el nombre se clasifican siempre como
\texttt{otro} para evitar que los documentos de anexos de pertinencia se confundan
con el documento principal. Los archivos dentro de subcarpetas denominadas
\texttt{old}, \texttt{antiguo}, \texttt{archivo} u \texttt{observaciones} se ignoran
completamente. Los archivos \texttt{.docx} que no caen en ninguna categoría y no
comienzan con \texttt{\textasciitilde\$} se tratan como programas de unidad de
enseñanza-aprendizaje (UEA).
\end{notabox}

\subsubsection{Salida}

Un objeto \texttt{InventarioLicenciatura} por licenciatura con las listas:
\texttt{pertinencia}, \texttt{plan\_estudios}, \texttt{mapa\_curricular},
\texttt{equivalencias}, \texttt{oficio}, \texttt{ueas}, \texttt{otros}.

%%──────────────────────────────────────────────────────────
\subsection{Agente~2 --- Extractor de plan (\texttt{extractor\_plan.py})}
%%──────────────────────────────────────────────────────────

\subsubsection{Función}

Extrae el contenido de los documentos de pertinencia y de plan de estudios
(archivos \texttt{.docx} o \texttt{.pdf}). Cumple dos propósitos:

\begin{itemize}[nosep]
  \item Proporcionar el texto completo del documento al Agente~3 (revisión con \IA{}).
  \item Extraer la tabla estructurada de UEAs del plan para la correspondencia (Agente~5).
\end{itemize}

\subsubsection{Extracción de texto}

Para archivos \texttt{.docx} utiliza \texttt{python-docx}: recorre todos los párrafos y
celdas de tabla y los concatena. Para archivos \texttt{.pdf} utiliza \texttt{pdfplumber}:
extrae el texto página a página.

\subsubsection{Extracción de la tabla de UEAs}

Busca dentro del documento tablas cuya primera fila contenga encabezados como
\texttt{clave}, \texttt{nombre de UEA} u \texttt{OBL./OPT.} y mapea las columnas por
nombre. Por cada fila posterior extrae: clave, nombre, tipo (OBL./OPT.),
horas teoría, horas práctica, créditos declarados y seriación.

\subsubsection{Salida}

Un objeto \texttt{DatosPlan} con los campos \texttt{texto} (cadena completa) y
\texttt{ueas} (lista de \texttt{RegistroUEA}).

%%──────────────────────────────────────────────────────────
\subsection{Agente~3 --- Revisor \IA{} (\texttt{revisor\_ia.py})}
%%──────────────────────────────────────────────────────────

\subsubsection{Función}

Realiza la revisión legislativa de fondo mediante la \textsc{api} de Claude
(\texttt{claude-sonnet-4-6}). Evalúa dos conjuntos de requisitos:

\begin{itemize}[nosep]
  \item \textbf{Sección~I}: el documento de pertinencia y justificación contra las
        fracciones I--X del Art.\ 34 \RES{}.
  \item \textbf{Sección~II}: el plan de estudios contra las fracciones I--XVI del
        Art.\ 44 \RES{}.
\end{itemize}

\subsubsection{Construcción del prompt}

Para cada sección se construye un prompt que incluye: (a)~el rol del revisor
(Consejo Divisional DCBI, UAM-Azcapotzalco), (b)~el texto del documento truncado a
12\,000 caracteres, (c)~la lista de fracciones a verificar con su texto normativo,
e (d)~instrucciones de evaluación. El modelo devuelve un arreglo JSON con una
entrada por fracción.

\subsubsection{Campos por fracción}

\begin{center}
\begin{tabular}{lp{9cm}}
\toprule
\textbf{Campo} & \textbf{Descripción} \\
\midrule
\texttt{cumple}      & Veredicto: \texttt{Sí} / \texttt{Parcial} / \texttt{No} / \texttt{N/A} \\
\texttt{evidencia}   & Cita textual o referencia de sección que respalda el veredicto
                       (máx.\ 120~car.) \\
\texttt{observacion} & Descripción de lo que falta o debe corregirse cuando el
                       veredicto es \texttt{No} o \texttt{Parcial} (máx.\ 160~car.) \\
\texttt{detalle}     & Análisis narrativo de una oración (máx.\ 180~car.) \\
\bottomrule
\end{tabular}
\end{center}

\begin{notabox}
La fracción~X del Art.\ 34 (residencias clínicas) está marcada como opcional:
el modelo puede asignarle \texttt{N/A} para las licenciaturas de ingeniería.
\end{notabox}

\subsubsection{Salida}

Una lista de objetos \texttt{ResultadoFraccion} con los campos \texttt{fraccion},
\texttt{fundamento}, \texttt{elemento}, \texttt{cumple}, \texttt{evidencia},
\texttt{observaciones} y \texttt{detalle}.

%%──────────────────────────────────────────────────────────
\subsection{Agente~4 --- Extractor de UEA (\texttt{extractor\_uea.py})}
%%──────────────────────────────────────────────────────────

\subsubsection{Función}

Procesa individualmente cada archivo de programa de UEA (\texttt{.docx}).
Extrae los datos del encabezado institucional y verifica la presencia de las secciones
requeridas por el Art.\ 45 \RES{}.

\subsubsection{Encabezado}

Lee la tabla de encabezado institucional (formato UAM) para obtener: clave de UEA,
nombre, tipo (OBL./OPT.), horas teoría, horas práctica, créditos declarados
y seriación.

\subsubsection{Verificación de secciones}

Busca en el texto completo del documento palabras clave asociadas a cada sección
requerida:

\begin{center}
\begin{tabular}{ll}
\toprule
\textbf{Sección} & \textbf{Palabras clave buscadas} \\
\midrule
Objetivos             & \texttt{objetivo}, \texttt{objetivos} \\
Contenido sintético   & \texttt{contenido}, \texttt{temario}, \texttt{contenido sintético} \\
Conducción            & \texttt{conducción}, \texttt{modalidades de conducción} \\
Evaluación            & \texttt{evaluación}, \texttt{modalidades de evaluación} \\
Bibliografía          & \texttt{bibliografía}, \texttt{referencias}, \texttt{bibliography} \\
\bottomrule
\end{tabular}
\end{center}

\subsubsection{Validación de créditos}

Calcula los créditos esperados según la fórmula del Art.\ 56 \RES{}:
\[
  \text{créditos} = 2 \times h_{\text{teoría}} + 1 \times h_{\text{práctica}}
\]
y los compara con el valor declarado en el encabezado del programa.

\subsubsection{Salida}

Un objeto \texttt{DatosUEA} por archivo, con todos los campos extraídos, los
indicadores booleanos de presencia de secciones, el campo \texttt{creditos\_ok}
y una lista de \texttt{advertencias} (texto libre para anomalías menores).

%%──────────────────────────────────────────────────────────
\subsection{Agente~5 --- Correspondencia (\texttt{correspondencia\_v2.py})}
%%──────────────────────────────────────────────────────────

\subsubsection{Función}

Cruza la lista de UEAs declaradas en el plan de estudios (extraída por el Agente~2)
con los archivos de programa procesados por el Agente~4.
Detecta tres tipos de problemas:

\begin{enumerate}[nosep]
  \item \textbf{Sin archivo en programas}: la UEA aparece en el plan pero no existe
        un archivo de programa correspondiente.
  \item \textbf{Sin UEA en plan}: existe un archivo de programa pero la UEA no
        aparece en la tabla del plan.
  \item \textbf{Discrepancias de campos}: la UEA existe en ambos lados pero hay
        diferencias en clave, nombre, tipo, horas o créditos.
\end{enumerate}

La correspondencia se establece por coincidencia de clave numérica y/o nombre
normalizado (minúsculas, sin tildes, sin signos de puntuación).

\subsubsection{Salida}

Una lista de objetos \texttt{ResultadoCorrespondencia} con los campos \texttt{clave},
\texttt{nombre\_plan}, \texttt{archivo}, \texttt{sin\_archivo}, \texttt{sin\_en\_plan},
\texttt{discrepancias} (lista de objetos \texttt{Discrepancia} con \texttt{campo},
\texttt{en\_plan} y \texttt{en\_programa}) y \texttt{ok}.

%%──────────────────────────────────────────────────────────
\subsection{Agente Excel (\texttt{excel\_v2.py})}
%%──────────────────────────────────────────────────────────

\subsubsection{Función}

Consolida todos los resultados y genera dos archivos Excel en la carpeta de salida
\texttt{Revisiones/}. Los archivos siguen la estructura del \textit{Formato de Revisión
de Planes y Programas} institucional y emplean la paleta de colores corporativa de
la UAM.

\subsubsection{Estructura de cada archivo}

Ambos archivos contienen una hoja de \textbf{RESUMEN} (índice cero del libro)
con la tabla consolidada de las diez licenciaturas, seguida de una hoja por
licenciatura con el detalle completo.

%%============================================================
\section{Documentos de salida}
%%============================================================

%%──────────────────────────────────────────────────────────
\subsection{Archivo 1: \texttt{Revision\_DCBI\_2026\_Planes\_<fecha>.xlsx}}
%%──────────────────────────────────────────────────────────

Contiene la revisión legislativa de los documentos principales (Secciones~I y~II).

\subsubsection{Hoja RESUMEN}

Tabla de siete columnas con una fila por licenciatura:

\begin{center}
\begin{tabular}{cl}
\toprule
\textbf{Col.} & \textbf{Contenido} \\
\midrule
A & Nombre de la licenciatura \\
B & Conteo de fracciones Art.\ 34 con veredicto \texttt{Sí} \\
C & Conteo de fracciones Art.\ 34 con veredicto \texttt{Parcial} \\
D & Conteo de fracciones Art.\ 34 con veredicto \texttt{No} \\
E & Conteo de fracciones Art.\ 44 con veredicto \texttt{Sí} \\
F & Conteo de fracciones Art.\ 44 con veredicto \texttt{Parcial} \\
G & Conteo de fracciones Art.\ 44 con veredicto \texttt{No} \\
\bottomrule
\end{tabular}
\end{center}

\subsubsection{Hojas por licenciatura}

Cada hoja contiene dos bloques con encabezado coloreado (azul UAM).

\paragraph{Bloque Sección~I --- Documentación General (Art.\ 34 \RES{})}

Una fila por cada fracción (I--X). Las diez fracciones evaluadas son:

\begin{center}
\renewcommand{\arraystretch}{1.2}
\begin{longtable}{cp{10.5cm}}
\toprule
\textbf{Fracc.} & \textbf{Elemento del Art.\ 34 \RES{}} \\
\midrule
\endhead
I    & Relevancia social y académica, pertinencia teórico-práctica, objetivos generales
       y estructura del plan; mapa curricular con perspectiva de género e inclusión;
       atención a necesidades y demandas de la sociedad; contribución al desarrollo
       científico, tecnológico, humanístico y cultural. \\[4pt]
II   & El alumnado a atender y el egreso previsible. \\[4pt]
III  & El perfil profesional de quienes egresen y su posible ocupación. \\[4pt]
IV   & La oferta de planes similares en otras instituciones de educación superior de
       la zona de influencia y la situación de sus egresados. \\[4pt]
V    & La población con prerrequisitos curriculares para demandar los estudios. \\[4pt]
VI   & Estimación de recursos necesarios: perfil docente y programa de formación;
       personal administrativo; infraestructura (instalaciones, laboratorios, acervo);
       impacto presupuestal. \\[4pt]
VII  & Las posibilidades de financiamiento, en su caso. \\[4pt]
VIII & Participación de órganos responsables si el plan es impartido por más de una
       división/unidad o en corresponsabilidad con otra institución. \\[4pt]
IX   & Información adicional a juicio del Consejo Divisional. \\[4pt]
X    & \textit{(N/A para ingenierías)} Instituciones para residencias clínicas. \\
\bottomrule
\end{longtable}
\end{center}

\paragraph{Bloque Sección~II --- Plan de Estudios Actualizado (Art.\ 44 \RES{})}

Una fila por cada fracción (I--XVI). Evalúa los dieciséis elementos mínimos
que debe contener el plan de estudios según el Art.\ 44, desde la denominación
de la unidad que imparte el plan (fracc.\ I) hasta la obligación de servicio social
(fracc.\ XVI).

\paragraph{Columnas de cada bloque (Secciones~I y~II)}

\begin{center}
\begin{tabular}{clp{7.5cm}}
\toprule
\textbf{Col.} & \textbf{Encabezado} & \textbf{Contenido} \\
\midrule
A & Fundamento          & Identificador de la fracción
                          (\textit{p.\,ej.}\ \texttt{Art.\ 34 fracc.\ III}) \\
B & Requisito           & Texto normativo del artículo \\
C & Cumple              & Veredicto: \texttt{Sí} / \texttt{Parcial} / \texttt{No}
                          / \texttt{N/A} (celda coloreada) \\
D & Evidencia en el doc.& Cita o referencia de sección que respalda el veredicto \\
E & Observación         & Descripción de lo que falta o debe corregirse \\
F & Análisis del revisor \IA{} & Análisis narrativo generado por el modelo \\
\bottomrule
\end{tabular}
\end{center}

\begin{okbox}
La columna~C se colorea según el veredicto: verde (\texttt{Sí}),
naranja (\texttt{Parcial}), rojo (\texttt{No}), gris (\texttt{N/A}).
\end{okbox}

%%──────────────────────────────────────────────────────────
\subsection{Archivo 2: \texttt{Revision\_DCBI\_2026\_Programas\_<fecha>.xlsx}}
%%──────────────────────────────────────────────────────────

Contiene la revisión de los programas individuales de UEA y la correspondencia
con el plan (Secciones~III y~IV).

\subsubsection{Hoja RESUMEN}

Tabla de once columnas con una fila por licenciatura:

\begin{center}
\begin{tabular}{cl}
\toprule
\textbf{Col.} & \textbf{Contenido} \\
\midrule
A & Nombre de la licenciatura \\
B & Total de UEAs procesadas \\
C & Conteo de UEAs con créditos correctos / total \\
D & Número de UEAs con error de créditos \\
E & UEAs sin sección de objetivos \\
F & UEAs sin contenido sintético \\
G & UEAs sin bibliografía \\
H & UEAs de correspondencia sin archivo de programa \\
I & Pares con discrepancias entre plan y programa \\
J & UEAs de correspondencia sin problema \\
K & Estado global (\texttt{OK} o ``$n$ obs.'', coloreado) \\
\bottomrule
\end{tabular}
\end{center}

\subsubsection{Hojas por licenciatura}

\paragraph{Bloque Sección~III --- Programas de Estudio (Art.\ 45 + Art.\ 56 \RES{})}

Una fila por UEA detectada. Columnas:

\begin{center}
\renewcommand{\arraystretch}{1.15}
\begin{tabular}{clp{7cm}}
\toprule
\textbf{Col.} & \textbf{Encabezado} & \textbf{Contenido} \\
\midrule
A & Clave UEA               & Clave numérica extraída del encabezado \\
B & Nombre de UEA           & Denominación \\
C & OBL/OPT                 & Tipo: obligatoria u optativa \\
D & Seriación               & Texto de seriación o «sin seriación» \\
E & Objetivos (Sí/No)       & Presencia de la sección de objetivos \\
F & Contenido (Sí/No)       & Presencia del contenido sintético \\
G & Conducción (Sí/No)      & Presencia de modalidades de conducción \\
H & Evaluación (Sí/No)      & Presencia de modalidades de evaluación \\
I & H.\ Teoría              & Horas de teoría declaradas \\
J & H.\ Prácticas           & Horas de práctica declaradas \\
K & Créditos declarados     & Créditos en el encabezado del programa \\
L & Créditos fórmula        & $2h_T + h_P$ (Art.\ 56 \RES{}) \\
M & Créditos OK             & Coincidencia K\,=\,L (verde/rojo) \\
N & Bibliografía (Sí/No)    & Presencia de sección de bibliografía \\
O & Advertencias            & Mensajes de error o anomalías menores \\
\bottomrule
\end{tabular}
\end{center}

\begin{notabox}
Las columnas E--H y N se colorean en verde cuando el valor es \texttt{Sí}
y en rojo cuando es \texttt{No}.
\end{notabox}

\paragraph{Bloque Sección~IV --- Correspondencia Plan vs.\ Programas de UEA}

Una fila por UEA cruzada. Columnas:

\begin{center}
\begin{tabular}{clp{7.5cm}}
\toprule
\textbf{Col.} & \textbf{Encabezado} & \textbf{Contenido} \\
\midrule
A & Clave (plan)              & Clave tal como aparece en el plan \\
B & Nombre en plan            & Denominación en la tabla del plan \\
C & Archivo programa          & Nombre del archivo \texttt{.docx} correspondiente \\
D & Sin archivo en programas  & \texttt{Sí} (rojo) si la UEA del plan carece de archivo \\
E & Sin UEA en plan           & \texttt{Sí} (rojo) si existe archivo pero no aparece en el plan \\
F & Discrepancias             & Lista en formato
                                \texttt{campo: valor\_plan} $\to$ \texttt{valor\_programa} \\
\bottomrule
\end{tabular}
\end{center}

\begin{notabox}
La columna~A (\textit{Clave}) se colorea en \textbf{verde} (\texttt{C6EFCE})
cuando la UEA no presenta ningún problema (\texttt{res.ok}\,=\,\texttt{True}).
Las columnas~D y~E se colorean en \textbf{rojo} (\texttt{FFC7CE}) con texto en negrita
cuando su valor es \texttt{Sí}.
La columna~F se colorea en \textbf{amarillo} (\texttt{FFEB9C}) cuando contiene
al menos una discrepancia.
\end{notabox}

%%============================================================
\section{Uso del sistema}
%%============================================================

%%──────────────────────────────────────────────────────────
\subsection{Ejecución}
%%──────────────────────────────────────────────────────────

\begin{lstlisting}[language={}]
cd revisor/

# Todas las licenciaturas (requiere ANTHROPIC_API_KEY)
python3 main_v2.py

# Solo una licenciatura
python3 main_v2.py -l Quimica

# Sin llamadas a la API de Claude (Secciones I y II en blanco)
python3 main_v2.py --no-ia

# Salida detallada por UEA
python3 main_v2.py -v
\end{lstlisting}

%%──────────────────────────────────────────────────────────
\subsection{Requisitos de entorno}
%%──────────────────────────────────────────────────────────

\begin{itemize}[nosep]
  \item Python 3.10 o superior.
  \item Paquetes: \texttt{anthropic}, \texttt{python-docx}, \texttt{pdfplumber},
        \texttt{openpyxl}.
  \item Variable de entorno \texttt{ANTHROPIC\_API\_KEY} o archivo
        \texttt{revisor/.env} con la clave.
\end{itemize}

%%──────────────────────────────────────────────────────────
\subsection{Carpeta de salida}
%%──────────────────────────────────────────────────────────

Los archivos Excel se escriben en:

\begin{lstlisting}[language={}]
Revisiones/
  Revision_DCBI_2026_Planes_<AAAA-MM-DD>.xlsx
  Revision_DCBI_2026_Programas_<AAAA-MM-DD>.xlsx
\end{lstlisting}

donde \texttt{<AAAA-MM-DD>} es la fecha del sistema en el momento de la ejecución.

%%============================================================
\section{Referencia legislativa}
%%============================================================

\begin{center}
\begin{tabular}{lll}
\toprule
\textbf{Artículo} & \textbf{Contenido} & \textbf{Sección en el Excel} \\
\midrule
Art.\ 34 \RES{} & Documentación General (pertinencia y justificación) & Sección~I \\
Art.\ 44 \RES{} & Contenido mínimo del plan de estudios               & Sección~II \\
Art.\ 45 \RES{} & Contenido mínimo de los programas de UEA            & Sección~III \\
Art.\ 56 \RES{} & Definición y fórmula de créditos                    & Sección~III \\
Art.\ 48 \RES{} & Límites de créditos de licenciatura (410--615)      & (resumen) \\
\bottomrule
\end{tabular}
\end{center}

%%============================================================
\end{document}
%%============================================================
